\subsection{The exact coefficient-five gate in a legal original factorial subfamily}

\begin{result}{E.199}{Proved}
{Minimal five-adic prime shells in the original A.44 recursion}
The original recursion E.102, with its unchanged exponents
\(n_j=(j+4)!\), has an infinite legal ordinary subfamily in which
\[
 \begin{gathered}
 s_j=v_5(n_{j+1}),\\
 q_j\equiv1+5^{s_j}\pmod{5^{s_j+1}},\qquad
 r_j\equiv1-5^{s_j}\pmod{5^{s_j+1}}.
 \end{gathered}
\tag{E.199.1}
\]
In particular \(v_5(q_j^2-1)=v_5(r_j^2-1)=s_j\).
This changes only permitted old-prime unit-root choices, and does not
scale the factorials or alter the radical splitting requirements.

For such an output take the unique possible global integral class-two
predecessor of B.101, with the literal coordinates of B.102, and put
\[
 \begin{aligned}
 D_{A,j}&=\gamma_j-q_j\alpha_j u_j,&
 d_{A,j}&=(q_j^2-1)/n_{j+1},\\
 D_{B,j}&=\widetilde\gamma_j+\beta_jv_j,&
 d_{B,j}&=(r_j^2-1)/n_{j+1}.
 \end{aligned}
\tag{E.199.2}
\]
Here the coefficient prime is five and both \(d_{A,j},d_{B,j}\) are
five-adic units.  The complete local coefficient-five third tests are
\[
       v_5(D_{A,j})\geq v_5(j+5),\qquad
       v_5(D_{B,j})\geq v_5(j+5)\quad(j\geq0).
\tag{E.199.3}
\]
Thus these local tests are automatic at every index not divisible by
five, including the second pair \(q_1,r_1\).  The first possible
nonautomatic tests are the two finite equations
\[
 \gamma_0-\alpha_0u_0=0\pmod5,\qquad
 \widetilde\gamma_0+\beta_0v_0=0\pmod5.
\tag{E.199.4}
\]
They are determined by a fixed legal prefix through index five.
No assertion that either equation is forced, or that they can always
be satisfied simultaneously, is made.

After E.199.3, the continuous coefficient-five third obstruction
vanishes if and only if
\[
 \frac{D_{A,j}}{j+5}\longrightarrow0,\qquad
 \frac{D_{B,j}}{j+5}\longrightarrow0
                    \quad\hbox{in }\mathbb Z_5.
\tag{E.199.5}
\]
The minimal shells therefore do not remove the continuous convergence
gate.  This statement isolates an exact subfamily of the original
problem; it neither constructs a class-three survivor nor excludes
all original A.44 outputs.
\end{result}

\begin{proof}
At stage \(j\), use the notation \(a=a_j,b=b_j,n=n_j,m=n_{j+1}\)
of E.102.  All selected support primes exceed five.  Thus \(a,b\)
are five-adic units and
\[
 a\equiv3\pmod5,\qquad b\equiv-2\pmod5,\qquad
 v_5(a+b-1)\geq H_j(5)=2v_5(m)+j+3.
\tag{E.199.6}
\]
The initial error is zero, so the same inequality holds initially.
Put \(c=v_5(n)\) and \(s=v_5(m)\); then \(1\leq s\) and
\(c\leq s\).  Choose the exact old local root coordinate
\(\xi=1+5^s\), and set
\[
                   Z=\frac{1-a\xi^n}{b}.
\tag{E.199.7}
\]
For a five-adic unit \(n/5^c\), the binomial expansion, or the
five-adic logarithm, gives
\[
 v_5(\xi^n-1)=c+s,\qquad
 v_5(Z-1)=c+s,
\tag{E.199.8}
\]
because the error in E.199.6 has strictly larger valuation.
The power map \(z\mapsto z^n\) sends \(1+5\mathbb Z_5\)
isomorphically onto \(1+5^{c+1}\mathbb Z_5\).
One may check this without a nonuniform Hensel division: the convergent
logarithm and exponential are mutually inverse isometries between
\(1+5\mathbb Z_5\) and \(5\mathbb Z_5\), and this power map is
multiplication by \(n\) on logarithms.  Since \(s\geq1\), there is
therefore a unique \(\eta\in1+5\mathbb Z_5\) with
\(\eta^n=Z\), and
\[
 v_5(\eta-1)=s,\qquad
 \frac{\eta-1}{5^s}\equiv
       -\frac ab\frac{\xi-1}{5^s}\equiv-1\pmod5.
\tag{E.199.9}
\]
This is an actual unit-coordinate solution of
\(a\xi^n+b\eta^n=1\).  Both coordinates are one modulo
\(5^{v_5(m)}\), exactly as required by the original recursion.
Its CRT precision is \(L_5=H_{j+1}(5)>s\), so the original
congruences for \(q_j\) and \(r_j\) retain precisely the two
residue classes in E.199.1.  At every other old place make any of
the original permitted choices.  The CRT unit classes remain one
modulo \(m\).  Consequently the same cyclotomic/radical intersection
proof and Chebotarev step of E.102 apply unchanged, including the
additional residue condition for \(r_j\) at \(q_j\) and the
radical exponent \(m\) at the second prime choice.
The original induction and its local limit proof therefore produce
an ordinary infinite-support output.  This argument works at every
stage, so it constructs the stated subfamily.

For each new support prime, \(q_j+1\) or \(r_j+1\) is a five-adic
unit.  E.199.1 gives the valuation assertion and the unit status of
\(d_{A,j},d_{B,j}\).  The exact normalized residues in B.102 are
\[
 T_{A,j}=\frac{D_{A,j}}{(j+5)d_{A,j}},\qquad
 T_{B,j}=\frac{D_{B,j}}{(j+5)d_{B,j}}.
\tag{E.199.10}
\]
All coordinates in the numerators are integral.  Coordinatewise
integrality is therefore exactly E.199.3, and is automatic when
\(5\nmid j+5\).  For \(j=0\), both primes are one modulo five,
so the two first congruences are exactly E.199.4.  At two and three
there is no coefficient-five local weight-three obstruction since
\(2^2-1\) and \(3^2-1\) are five-adic units.  At the coefficient
place five, finite local duality gives zero because the mod-five
cyclotomic character squared is nontrivial; the real group has no
five-primary positive cohomology.  At all other primes the zero
valuation residue calculation in B.102 applies.  This accounts for
all local places in the coefficient-five component.

We make the finite-prefix assertion precise.  For B.20 with
\(\ell=5,N=M=1\), the first factorial divisible by five occurs
at \(J_M=1\).  The relevant old primes are \(2,3,q_0,r_0\),
and the largest \(v_5(p-1)\) among them is one.  Thus the division
budget there is \(D=2\).  The first \(n_L\) divisible by twenty-five
is \(n_6=10!\), so fixing the prime prefix through index five
fixes the reduction modulo five of the unique integral class-two
class.  The literal abelian cochains are fixed modulo twenty-five.
At \(q_0,r_0\), division by \(p-1\) loses exactly one five-adic
digit; hence \(u_0,v_0\pmod5\) are fixed too.  These are precisely
the finite inputs in E.199.4.  Central or aligned abelian changes
of representative do not change the resulting cohomological tests.

Finally, multiplication by arbitrary five-adic units preserves
convergence to zero.  The exact continuous localization theorem
A.22 and B.102 therefore turn E.199.10 into E.199.5, with no
replacement of the actual image by its closure.  In particular
coordinatewise integrality at the sparse exceptional indices does
not imply E.199.5.  The factorial in the numerator has been cancelled
by the equally growing five-primary part of the denominator, so it
cannot be cited a second time as a proof of that convergence.
\end{proof}
